% answers_short.tex —— 2005 年·普通高等学校招生考试·全国卷Ⅲ·理科 · 简短答案速查
% 由 tools/make_short_answers.py 自 answers.tex 的 \gksolhead/\gkanswer 提取，
% 并与 2005/2005_数学真题_全国卷Ⅲ理科_答案要点.md 逐题比对后生成，勿手改（改源 MD 后重跑）。

\begin{gknotice}
\textbf{使用说明}：本节只给最终答案便于快速核对，完整推导见解析版。答案为编写者独立推导，\textbf{非官方标准答案}。
\end{gknotice}

\gkpart{选择题}
\begin{center}\begin{tabular}{ccccccccccccc}
\toprule
题号 & 1 & 2 & 3 & 4 & 5 & 6 & 7 & 8 & 9 & 10 & 11 & 12 \\
\midrule
答案 & D & B & B & C & A & C & C & B & C & D & D & A \\
\bottomrule
\end{tabular}\end{center}

\gkpart{填空题}
\begin{itemize}[leftmargin=2.6em]
  \item[\textbf{13.}] $z = 1 - \dfrac{3}{2}\mathrm{i}$。
  \item[\textbf{14.}] $k = -\dfrac{2}{3}$。
  \item[\textbf{15.}] $E\xi = \dfrac{4}{7}$。
  \item[\textbf{16.}] $3$。
\end{itemize}

\gkpart{解答题}
\begin{itemize}[leftmargin=2.6em]
  \item[\textbf{17.}] (1) $P(\text{\text{甲}}) = 0.2$，$P(\text{\text{乙}}) = 0.25$，$P(\text{\text{丙}}) = 0.5$；(2) $0.7$
  \item[\textbf{18.}] (1) 由平面 $VAD \perp$ 底面及 $AB \perp AD$ 得 $AB \perp$ 平面 $VAD$；(2) $\arccos\dfrac{\sqrt{21}}{7}$
  \item[\textbf{19.}] (1) $\cot A + \cot C = \dfrac{4\sqrt{7}}{7}$；(2) $a + c = 3$
  \item[\textbf{20.}] $k_n = 3^{n+1}$。
  \item[\textbf{21.}] (1) 当且仅当 $x_1 + x_2 = 0$ 时直线 $l$ 经过焦点 $F$；(2) $\left[\dfrac{9}{32}, +\infty\right)$
  \item[\textbf{22.}] (1) $f(x)$ 在 $\left[0, \dfrac{1}{2}\right]$ 上单调递减、在 $\left[\dfrac{1}{2}, 1\right]$ 上单调递增，值域 $\left[-4, -3\right]$；(2) $1 \leqslant a \leqslant \dfrac{3}{2}$
\end{itemize}
